\section{Uniform spectral rates and phase products}
\label{S:uniform-physical-products}

Retain every original covariance and amplitude.
Set $E=H^{1/2}\times H^{-1/2}$, $F=E\oplus\R^2$,
and $m_R=\eta(\cdot/R)^2$, $R\ge1$, with the same fixed root profile.
Choose
\[
  \begin{aligned}
    &\frac12<\theta<\min(1-a/2,(3-a-\alpha)/2),\qquad \nu=1/\theta,\\
    &(1-\theta)^{-1}<r<2/a,\qquad \nu\le p<2<r .
  \end{aligned}
\]
Choose auxiliary exponents
$0<\gamma<\min(1/2,(2-\alpha)/2)$ and
$\theta<\beta<\min(1-a/2,(3-a-\alpha)/2)$ for the actual phase construction.
The original strict threshold makes these intervals nonempty.
Write $i_I:H_I\to H_W$ for the supported closed inclusion.
For nonzero covariance the actual field is
\[
  \mathcal T_R=T_R(C_t\otimes I_{H_{\rm sp}}),\qquad
  \|\mathcal T_R\|_{H_I\to L^r(I;\mathcal K(E))}
  \le C_{\rm mean}=M_0\|C_t\| ,
\]
by Proposition~\ref{P:uniform-mean-regulator}.
For zero covariance set $H_{\rm sp}=H_I=\{0\}$ and
$M_0=C_{\rm mean}=0$. Means and their rates are then zero,
and mean flows are free; no nonzero-noise estimate is needed for that case.
The residual comparisons below still act on the nonzero phase $E$.

\subsection{The actual E class and its own spectral constants}

For real smooth compact $q$ let $V_q=m_R(\Gamma*q)$,
$M_q:H^{1/2}\to H^{-1/2}$ be its multiplier and
$L_q^E(u,v)=(0,M_qu)$.
Let $\mathcal V_{E,R}$ be their operator-norm closure in $\mathcal K(E)$.

\begin{lemma}[Unique potential and fixed-radius lift]\label{L:uniform-E-lift}
  Each $L\in\mathcal V_{E,R}$ has the form $L(u,v)=(0,Mu)$,
  with $\|L\|=\|M\|=:k$, and determines a unique real
  $V\in H^{-1/2}$ supported in $\operatorname{supp}m_R$.
  The actual symmetric pairing satisfies
  \begin{equation}\label{E:uniform-E-pairing}
    |v_V(u,w)|\le k\|u\|_{H^{1/2}}\|w\|_{H^{1/2}} .
  \end{equation}
  It has a lift $L_F$ in the exact class of
  Section~\ref{S:signed-static-wave}, with $L_Fi=iL$.
  Lift constants may depend on $R$ and are used only for identification.
  Here $v_V(u,w)=\langle Vu,w\rangle_{H^{-1/2},H^{1/2}}$ on real
  states, with its Hermitian extension on complex states.
\end{lemma}

\begin{proof}
  For zero covariance, $L=M=V=L_F=0$ gives the claims directly.
  If $L_n^E\to L$, the zero columns persist and $M_n\to M$
  in $\mathcal L(H^{1/2},H^{-1/2})$, with the stated norm equality.
  At fixed $R$ choose real compact smooth $\chi_R=1$ near
  $\operatorname{supp}m_R$.
  Since $M_n\chi_R=V_n$,
  $\|V_n-V_j\|_{H^{-1/2}}\le\|\chi_R\|_{H^{1/2}}\|M_n-M_j\|$.
  Thus, $V_n\to V$, with the same support by testing outside it.
  Smooth multiplication on $H^{-1/2}$ identifies $Mu=Vu$
  on compact tests, then on $H^{1/2}$ by density.
  The identity $M\chi_R=V$ gives uniqueness; symmetry and
  \eqref{E:uniform-E-pairing} pass to the limit.

  The original $F$ operators are
  $(u,v,c_1,c_2)\mapsto(0,M_nu+c_1V_n,0,0)$.
  Their differences are bounded by
  $(1+\|\chi_R\|_{H^{1/2}})\|M_n-M_j\|$.
  They therefore converge in the old $F$ class and give the lift.
  This bound is allowed to grow with $R$; it will enter no uniform estimate.
  Dense smooth raw controls and bounded completion of $T_R$
  also show that all completed mean fields lie in
  $\mathcal V_{E,R}$ almost everywhere.
  Bochner averages preserve this closed linear class.
  Complexification preserves $k$ by the Hilbert real/imaginary splitting.
  The real symmetric pairing becomes Hermitian and its diagonal
  splits into the two real diagonals.
  This proves Lemma~\ref{L:uniform-E-lift}.
\end{proof}

\begin{proposition}[Spectral rate in E]
  \label{P:uniform-E-spectral}
  For $L\in\mathcal V_{E,R}$ put
  \[
    \beta(V)=\inf_{\substack{u\in H^1\\\|u\|_2=1}}
    \{\|\nabla u\|_2^2-v_V(u,u)\},\qquad \lambda(L)=\sqrt{-\beta(V)} .
  \]
  If $K\ge\max(1,k)$, the actual physical wave satisfies, for $t\ge0$,
  \begin{equation}\label{E:uniform-E-static-bound}
    -K^2\le\beta(V)=\inf\operatorname{spec}H_V\le0,\qquad
    \|P_V^E(t)\|\le8(1+K)^2(1+t)e^{\lambda(L)t}.
  \end{equation}
  If two such operators have norms at most $K$, then
  \begin{equation}\label{E:uniform-E-rate-continuity}
    |\beta(V)-\beta(W)|\le(K+2)\|L_V-L_W\|,\qquad
    |\lambda(L_V)-\lambda(L_W)|
    \le\sqrt{(K+2)\|L_V-L_W\|}.
  \end{equation}
  Every quantitative constant here uses the $E$ norm alone.
\end{proposition}

\begin{proof}
  We derive the budgets without transferring a quantitative $F$ bound.
  With $b=\|u\|_2$, $z=\|u\|_{H^1}$, Fourier Cauchy--Schwarz gives
  $\|u\|_{H^{1/2}}^2\le bz$.
  Equation~\eqref{E:uniform-E-pairing} and the half-weight
  arithmetic--geometric mean inequality give
  $|v_V(u,u)|\le Kbz\le z^2/2+K^2b^2/2$.
  The weights are $1/2,1/2$, bases $z^2,(Kb)^2\ge0$;
  positive powers include zero norms.
  For $q_V=\|\nabla u\|_2^2-v_V(u,u)$ and $c=2+K^2$,
  \begin{equation}\label{E:uniform-E-form-budget}
    \tfrac12z^2+(1+K^2/2)b^2
    \le q_V(u)+cb^2
    \le\tfrac32z^2+(1+3K^2/2)b^2
    \le3(1+K^2)z^2 .
  \end{equation}
  This is a densely defined closed semibounded Hermitian form
  on full $H^1$: its shifted norm is equivalent to the complete $H^1$ norm.
  The lower bound dominates $2b^2$, since $z\ge b$ and $K\ge1$.
  From~\cite{teschl:14:mathematical}*{Theorem~2.14, p.~81, (2.63)}
  and Theorem~3.6, p.~109, (3.51)--(3.53), it represents a
  unique self-adjoint $H_V$ with $A=H_V+c\ge2$ and $D(A^{1/2})=H^1$.
  Closedness uses the definition on that book's p.~80.
  The spectral meaning of $A\ge2$ follows from bounded spectral slices:
  a nonzero slice below two would give a unit vector in $D(A)$
  contradicting the form lower bound. The bottom argument below
  uses the same representation.

  Put $B_0=1-\Delta$ and $L_K=3(1+K^2)$.
  The common forms obey $\tfrac12B_0\le A\le L_KB_0$.
  For positive self-adjoint $T,S\ge1$ with ordered forms,
  completing the square gives
  \[
    \langle f,(T+t)^{-1}f\rangle
    =\sup_u\{2\operatorname{Re}\langle f,u\rangle-q_T(u)-t\|u\|^2\}
    \ge\langle f,(S+t)^{-1}f\rangle .
  \]
  The identity
  $\sqrt x=\pi^{-1}\int_0^\infty t^{-1/2}x/(x+t)\ud t$
  follows by $t=xv^2$.
  Positive spectral integration and Tonelli give square-root
  form order, with domain inclusions, from
  $I-t(T+t)^{-1}\le I-t(S+t)^{-1}$.
  Apply this to $(B_0,2A)$ and $(A,L_KB_0)$, both pairs at least one:
  \[
    D(A^{1/4})=H^{1/2},\qquad
    2^{-1/4}\|u\|_{H^{1/2}}\le\|A^{1/4}u\|_2
    \le L_K^{1/4}\|u\|_{H^{1/2}} .
  \]
  Under the $L^2$ pairing the dual norm is $\|A^{-1/4}v\|_2$,
  by Cauchy--Schwarz and the equality test $u=A^{-1/2}v$.
  Completion gives $H^{-1/2}$.
  Thus, with $M_K=L_K^{1/4}$, the phase norm
  $\|(u,v)\|_{E_A}^2=\|A^{1/4}u\|_2^2+\|A^{-1/4}v\|_2^2$
  lies between $M_K^{-2}\|(u,v)\|_E^2$ and $M_K^2\|(u,v)\|_E^2$.
  The condition cost is $M_K^2=\sqrt{3(1+K^2)}\le2(1+K)$.

  For unit $L^2$ data, $b=1$ and $z\ge1$.
  The form budget gives $q_V\ge z^2/2-1-K^2/2\ge-K^2$.
  A dilated compact unit bump translated wholly off the potential
  support has zero potential pairing and kinetic energy tending to zero,
  so $\beta\le0$.
  Complex form infima agree with real ones: dividing each complex
  diagonal by its squared norm is the convex combination of the
  two normalized real diagonals, omitting zero parts.
  The spectral form formulas identify $\beta$ with the spectral
  bottom by bounded spectral slices.
  There are support points because scalar spectral measures of unit
  vectors have mass one; their arbitrarily small bounded slices
  give approximate eigenvectors and the reverse form bound.
  If $\Sigma$ denotes the closed PVM support and real $z\notin\Sigma$,
  the Borel function $f_z(h)=(h-z)^{-1}$ on $\Sigma$, zero outside,
  is bounded, as is $hf_z(h)$.
  Its calculus has range in $D(H_V)$ and supplies both inverse identities
  for $H_V-z$ on their respective domains. Thus no such $z$ is spectral;
  the approximate eigenvectors give the other inclusion.
  Hence $\lambda\le K$ and $h+c\ge2$ on the spectrum.

  Define $\operatorname{Cos}_t(h),\operatorname{Sin}_t(h)$ by
  $\cos(t\sqrt h),\sin(t\sqrt h)/\sqrt h$ for $h>0$,
  by $1,t$ at zero and by
  $\cosh(t\sqrt{-h}),\sinh(t\sqrt{-h})/\sqrt{-h}$ for $h<0$.
  In the $E_A$ isometry the matrix is
  \[
    \begin{pmatrix}
      \operatorname{Cos}_t(h)&\sqrt{h+c}\operatorname{Sin}_t(h)\\
      -h\operatorname{Sin}_t(h)/\sqrt{h+c}&\operatorname{Cos}_t(h)
    \end{pmatrix}.
  \]
  For $h\ge0$, its off-diagonals have bounds $1+\sqrt c\,t$ and one;
  at zero their values are $\sqrt c\,t,0$.
  For $\beta\le h<0$, the entries are bounded using
  $|\operatorname{Cos}_t|\le e^{\lambda t}$,
  $|\operatorname{Sin}_t|\le te^{\lambda t}$,
  upper off-diagonal at most $\sqrt c\,te^{\lambda t}$,
  and lower at most $Ke^{\lambda t}$.
  Since $\sqrt c\le2K$, each entry is at most
  $2(1+K)(1+t)e^{\lambda t}$.
  Set these entries to zero off the closed spectral support
  before applying bounded Borel calculus.
  From~\cite{teschl:14:mathematical}*{Theorems~3.1--3.2,
  pp.~102--104, (3.18)--(3.33)}, their matrix norm is at most
  twice that entry bound, and the condition cost gives
  \eqref{E:uniform-E-static-bound}.
  Spectral dominated convergence and scalar addition give the group.

  Denote the spectral matrix evolution by $P_V^E(t)$.
  It is the actual physical $E$ wave.
  The fixed-$R$ lift in Lemma~\ref{L:uniform-E-lift} belongs
  to the precise class in Proposition~\ref{P:signed-physical-wave}.
  Its form is this same $q_V$ and has the same self-adjoint
  representative by uniqueness, regardless of the lift's larger shift.
  On dense bounded spectral cutoffs the two waves have the same
  $\operatorname{Cos},\operatorname{Sin}$ functions of $H_V$.
  Boundedness extends equality to all $E$.
  The old graph-core Duhamel identification therefore applies;
  no lift norm or $\chi_R$ constant entered the budgets above.
  Conjugation commutes with the real structures, so real restrictions
  are the actual phases.
  For compact $L$, finite-rank approximation and strong continuity of
  the free group and its adjoint make
  $v\mapsto S_E(-v)LS_E(v)$ norm continuous on finite intervals.
  Its primitive is Lipschitz. Its Young equation is therefore the
  ordinary bounded-interaction equation by
  Proposition~\ref{P:constant-young-linear}.
  This identifies the averaged static baseline without transferring
  any lift norm into the uniform budgets.

  Finally, take a unit $\varepsilon$-minimizer for $W$,
  $0<\varepsilon\le1$. Since $\beta(W)\le0$, it has $q_W\le1$.
  Thus $z^2-1-Kz\le1$, which forces $z\le K+2$.
  The multiplier difference changes its form by at most
  $\|L_V-L_W\|\|u\|_{H^{1/2}}^2\le(K+2)\|L_V-L_W\|$.
  Take infima, let $\varepsilon\downarrow0$ and reverse $V,W$.
  The square-root inequality on nonnegative $-\beta$ gives
  \eqref{E:uniform-E-rate-continuity}.
  For zero covariance the class is $\{0\}$, $H_V=-\Delta$ and
  $\beta=\lambda=0$; the free wave gives the claims directly.
  This proves Proposition~\ref{P:uniform-E-spectral}.
\end{proof}

\subsection{Scalar-time averaging and actual mean budgets}

\begin{proposition}[Uniform completed-field averaging]\label{P:uniform-rate-averaging}
  Use the true slow partition $I_m$ with mesh at most $\delta B$.
  Let $\mathcal A_\delta$ average over each positive-length cell;
  an empty final cell is ignored.
  Put $\omega_\delta=\|(\mathcal A_\delta-I)C_t\|_{H_t\to L^r(I)}$.
  At fixed $B$, $\omega_\delta\to0$.
  For $M\ge0$, if $\|\Phi\|_{H_I}\le M$, $W_R=\mathcal T_R\Phi$,
  $\overline W_R=\mathcal A_\delta W_R$ and $Z=C_{\rm mean}M$, then
  \begin{align}
    \|W_R\|_{L^r},\ \|\overline W_R\|_{L^r}&\le Z,\qquad
    \|\overline W_R-W_R\|_{L^r}\le M_0M\omega_\delta=:\Omega_\delta(M),
    \label{E:uniform-field-averaging}\\
    \operatorname{Rate}_R(\Phi)&=\int_I\lambda(\mathcal T_R\Phi(s))\ud s,
    \label{E:uniform-mean-rate}\\
    \int_I|\lambda(\overline W_R)-\lambda(W_R)|\ud s
    &\le\sqrt{3+2Z}\sqrt{\Omega_\delta(M)}=:\Lambda_\delta(M)\to0 .
    \label{E:uniform-rate-averaging}
  \end{align}
  These errors are uniform in $R\ge1$.
  The rate is finite, nonnegative and continuous on $H_I$.
\end{proposition}

\begin{proof}
  Bochner averaging is a contraction on $L^r(I;\mathscr Z)$
  for any Banach $\mathscr Z$, by Jensen.
  For continuous fields on $[0,1]$, its error is bounded by
  the modulus of continuity at the mesh.
  Simple-field approximation, followed by continuous approximation
  of scalar indicators in $L^r$, proves density of those fields.
  The contraction gives strong convergence for every Bochner field.
  In the scalar case it is uniform on the compact closure of
  $C_t$'s unit-ball image, by a finite net and
  $\|\mathcal A_\delta-I\|\le2$. This proves $\omega_\delta\to0$.

  On finite tensors, integration commutes with the fixed bounded $T_R$:
  \[
    \mathcal A_\delta^{\mathcal K(E)}
    T_R(C_t\otimes I)\Phi
    =T_R((\mathcal A_\delta C_t)\otimes I)\Phi .
  \]
  The dimension-free Hilbert tensor bound in
  Lemma~\ref{L:uniform-mean-tensor}, and $\|T_R\|\le M_0$,
  extend this typed identity to the completed $H_I$.
  Applying it to the difference gives~\eqref{E:uniform-field-averaging}.
  Averages remain in the same closed $E$ operator class.
  Only scalar-time compactness was used; the tensor map is not
  claimed compact. Nor does this assert noncompactness of a composite
  spatial family: each fixed-radius completed mean is compact.

  For any two such fields, apply~\eqref{E:uniform-E-rate-continuity}
  with the pointwise bound
  \[
    K(s)=\max(1,\|W_1(s)\|,\|W_2(s)\|).
  \]
  Cauchy--Schwarz gives
  \[
    \int_I|\lambda(W_1)-\lambda(W_2)|
    \le\sqrt{3+\|W_1\|_{L^1}+\|W_2\|_{L^1}}\,
    \sqrt{\|W_1-W_2\|_{L^1}} .
  \]
  Continuity of $\lambda$ on the class makes these quantities measurable.
  Since $|I|=1$, this proves~\eqref{E:uniform-rate-averaging}.
  It also proves rate continuity, while
  $\lambda(L)\le\max(1,\|L\|)$ gives
  $\operatorname{Rate}_R(\Phi)\le1+C_{\rm mean}\|\Phi\|$.
  For zero covariance the fields and rate errors are zero directly.
  This proves Proposition~\ref{P:uniform-rate-averaging}.
\end{proof}

For the remaining proofs use the full parameter package and true
$B,\delta,N,\ell_m,I_m$ partition of Section~\ref{S:actual-cameron-mean}.
For $W_R=\mathcal T_R\Phi$ define
\[
  H_m(u)=\int_0^{\min(u,\ell_m)}
  S_E(-v)W_R(s_m+\delta v)S_E(v)\ud v,\qquad
  P_m^{\rm mean}=S_E(\ell_m)U_{H_m}(\ell_m,0).
\]

\begin{lemma}[Explicit physical primitive and growth budgets]
  \label{L:uniform-product-budgets}
  Write $G_B=\sup_{|v|\le B}\|S_E(v)\|$ and
  $C_B=G_B^2[B^{1-1/r}+B^{1-1/r-\theta}]$.
  Then $1\le G_B\le3B$, $C_B\le18B^{3-1/r}\le32B^{3-1/r}$
  and $C_BB^{1/r}\ge B$.
  The stopped primitives lie in the zero-start little H\"older space and
  \begin{equation}\label{E:uniform-product-primitive}
    \|H_m\|_\theta\le a_m:=C_B\delta^{-1/r}\|W_R\|_{L^r(I_m)},
    \qquad\delta\sum_m a_m^r\le C_B^r\|W_R\|_{L^r}^r .
  \end{equation}
  There is a fixed $C_{\rm gr}=C_{\rm gr}(\theta)>0$ such that
  \begin{equation}\label{E:uniform-product-growth}
    \log\|S_E(\ell)U_D(\ell,0)\|
    \le C_{\rm gr}B(1+\|D\|_\theta^\nu),\qquad 0\le\ell\le B .
  \end{equation}
  Here $D$ is any zero-start $C^\theta([0,B];\mathcal L(E))$ path;
  no Gaussian or Hamiltonian hypothesis is imposed on this growth bound.
  The actual mean endpoints have $\|P_m^{\rm mean}\|\ge1$.
\end{lemma}

\begin{proof}
  The balanced free matrix has upper entry
  $\langle k\rangle\sin(vk)/k$ and lower $-k\sin(vk)/\langle k\rangle$.
  Splitting $k\le1$ and $k\ge1$ bounds the upper by $\sqrt2 B$,
  while the lower and diagonals have bound one.
  At zero the off-diagonals are $v,0$.
  The Frobenius bound gives $G_B\le\sqrt{3+2B^2}\le3B$;
  $S_E(0)=I$ gives $G_B\ge1$.
  The stated $C_B$ inequalities follow, including
  $C_BB^{1/r}=G_B^2(B+B^{1-\theta})\ge B$.
  H\"older in time with $1-1/r>\theta$, the affine time change,
  compact conjugation and constant stopping prove
  \eqref{E:uniform-product-primitive}; the true partition sums
  the $r$th powers, with zero on an empty cell.

  For the explicit Young budget put $d=\|D\|_\theta$ and
  $c_\theta=2\max(1,C_{\rm Young})$,
  as in the local estimate proving Lemma~\ref{L:young-upper-growth}.
  If $c_\theta dB^\theta\le1/4$, whole-interval contraction gives
  $\log\sup\|U_D\|\le\log(4/3)<1/2$.
  Otherwise at most $2B(4c_\theta d)^\nu$ contraction pieces
  have norm at most $4/3$.
  Combine their logarithms with $\log G_B\le\log(4B)$.
  For example
  $C_{\rm gr}=4+\log4+2(4c_\theta)^\nu$ gives
  \eqref{E:uniform-product-growth}, since $B\ge1$.
  This is the generic algebra argument on $\mathcal L(E)$
  with multiplication norm one, not an unspecified radius-dependent bound.

  Finally Proposition~\ref{P:symplectic-block-mean} applies exactly:
  $\mathcal T_R$ is the same compression of the actual $F$ mean,
  with the same one-mask physical model.
  Its fast direction is
  $\delta^{-a/2}T_{\delta,m}^*i_I\Phi$, finite at fixed $\delta$.
  The completed covariance identity identifies the primitive up to
  $\ell_m$, then stops it.
  Thus, pure-mean symplectic conservation gives the endpoint norm
  at least one, deterministically for each $\Phi$.
  No fast-direction cost estimate is used.
  An empty endpoint is $I$, and the zero-covariance case is the
  free identity in Lemma~\ref{L:symplectic-free}.
  This proves Lemma~\ref{L:uniform-product-budgets}.
\end{proof}

\subsection{Full blocks, the true last cell and polynomial control balls}

\begin{proposition}[Finite-mesh physical mean-product bound]
  \label{P:uniform-mean-product}
  Fix $B$ as above, $M\ge0$ and $K\ge1$, and put $Z=C_{\rm mean}M$.
  Uniformly for $R\ge1$, $\|\Phi\|\le M$,
  \begin{align}
    \delta\sum_{m=0}^N\log\|P_m^{\rm mean}\|
    -\operatorname{Rate}_R(\Phi)
    \le{}&\frac{L_{B,K}}B+
    C_{\rm gr}32^r Z^rB^{3r}(K^{-r}+K^{\nu-r})\notag\\
    &+\Lambda_\delta(M)+C(\theta,B,K)E_\delta+
    \operatorname{Last}_\delta,\label{E:uniform-mean-product}
  \end{align}
  where $L_{B,K}=\log16+4\log(1+K)+3\log(1+B)$ and
  \begin{align*}
    E_\delta&=(2/B)^{1-1/r}C_B\Omega_\delta(M)
    +(2/B)^{1-\nu/r}C_B^\nu\Omega_\delta(M)^\nu,\\
    \operatorname{Last}_\delta&=
    C_{\rm gr}B[\delta+C_B^\nu Z^\nu\delta^{1-\nu/r}].
  \end{align*}
  The last three terms vanish as $\delta\downarrow0$ at fixed $B,K,M$.
\end{proposition}

\begin{proof}
  On each positive-length cell use its own Bochner average $\overline M_m$;
  set it to zero on an empty cell and form the constant-field primitive
  $\overline H_m$.
  Put $e_m=C_B\delta^{-1/r}
  \|W_R-\overline W_R\|_{L^r(I_m)}$.
  Jensen and~\eqref{E:uniform-product-primitive} give
  \[
    \|\overline H_m\|_\theta\le a_m,\quad
    \|H_m-\overline H_m\|_\theta\le e_m,\quad
    \delta\sum e_m^r\le C_B^r\Omega_\delta(M)^r,\quad
    \delta\sum a_m^r\le D_M:=C_B^rZ^r .
  \]
  Finite power mean at exponents $r$ and $r/\nu>1$,
  with nonnegative inputs $e_m,e_m^\nu$, gives
  $\delta\sum_m(e_m+e_m^\nu)\le E_\delta$.
  Here $n=N+1\ge2$, $\delta>0$ and $\delta n\le2/B$;
  all subsequent powers are positive and zero errors are allowed.

  On full cells only,
  $\|\overline M_m\|\le(\delta B)^{-1/r}\|W_R\|_{L^r(I_m)}
  =a_m/(C_BB^{1/r})\le a_m/B$.
  Split full indices $m<N$ into good $a_m\le K$ and bad $a_m>K$.
  On a good block, Proposition~\ref{P:young-upper-multiplicative}
  applies with base $\overline H_m$, perturbation $H_m-\overline H_m$
  and the actual left multiplier $S_E(B)$.
  The baseline is the actual static wave of $\overline M_m$
  by Proposition~\ref{P:uniform-E-spectral}.
  Here $\|\overline M_m\|\le K/B\le K$, whereas
  $\|\overline H_m\|_\theta\le a_m\le K$.
  The spectral estimate gives
  \[
    \|S_E(B)U_{\overline H_m}(B,0)\|
    \le8(1+K)^2(1+B)e^{B\lambda(\overline M_m)}.
  \]
  Consequently,
  \[
    \log\|P_m^{\rm mean}\|
    \le L_{B,K}+B\lambda(\overline M_m)
    +C(\theta,B,K)(e_m+e_m^\nu).
  \]
  This sufficient prefactor follows from the bound in $E$,
  $8(1+K)^2(1+B)$, by enlarging it.
  Since $\delta N\le1/B$, the prefactors sum to at most $L_{B,K}/B$.
  Nonnegativity of $\lambda$ gives
  $\delta\sum_{\rm good}B\lambda(\overline M_m)
  \le\int_I\lambda(\overline W_R)
  \le\operatorname{Rate}_R(\Phi)+\Lambda_\delta(M)$.
  Bad and last rate cells can be dropped in this upper bound.

  On bad full cells~\eqref{E:uniform-product-growth} gives
  $C_{\rm gr}B(1+a_m^\nu)$.
  Because $a_m>K\ge1$ and $r>\nu$,
  $1+a_m^\nu\le(K^{-r}+K^{\nu-r})a_m^r$.
  Their weighted sum is bounded by
  $C_{\rm gr}BC_B^rZ^r(K^{-r}+K^{\nu-r})$.
  The inequality $C_B\le32B^{3-1/r}$ yields the displayed bad-block term.

  If $\ell_N=0$, the last endpoint is $I$ and
  $\log\|I\|=0$.
  Otherwise use the actual stopped path, not its frozen full-cell formula.
  Since $a_N\le C_B\delta^{-1/r}Z$, its weighted logarithm is at most
  $\operatorname{Last}_\delta$ by~\eqref{E:uniform-product-growth}.
  This tends to zero because $\nu<r$, even for arbitrarily short
  last cells; no division by that cell length occurs.
  Combining these estimates proves~\eqref{E:uniform-mean-product}.
  All constants are independent of $R$.
  This proves Proposition~\ref{P:uniform-mean-product}.
\end{proof}

\begin{corollary}[Every polynomial control radius]\label{C:uniform-polynomial-balls}
  For $c_{\rm ball}>0$, $\kappa_{\rm ball}\ge0$ put
  $M_B=c_{\rm ball}B^{\kappa_{\rm ball}}$.
  Then
  \[
    \limsup_{\substack{B\to\infty\\B\in\mathbb N}}
    \limsup_{\delta\downarrow0}
    \sup_{\substack{R\ge1\\\|\Phi\|\le M_B}}
    \left\{\delta\sum_m\log\|P_m^{\rm mean}\|
    -\operatorname{Rate}_R(\Phi)\right\}\le0 .
  \]
\end{corollary}

\begin{proof}
  Choose $s>r(3+\kappa_{\rm ball})/(r-\nu)$ and $K_B=B^s$.
  At each fixed $B$, the inner $\delta$ limit has a fixed finite
  radius and threshold, so the last three errors in
  \eqref{E:uniform-mean-product} vanish first.
  Its bad terms become
  \[
    C_{\rm gr}32^r C_{\rm mean}^r c_{\rm ball}^r
    \{B^{r(3+\kappa_{\rm ball})-sr}
    +B^{r(3+\kappa_{\rm ball})+s(\nu-r)}\}.
  \]
  Both exponents are negative: the second by choice and the first
  because $s>3+\kappa_{\rm ball}$.
  Also
  $L_{B,K_B}/B\le[\log16+7\log2+(4s+3)\log B]/B\to0$.
  Thus, the outer integer-$B$ limit gives the claim.
  No diagonal choice $B_\delta,M_\delta$ is asserted.
  At fixed parameters all primitive maps are bounded linear into
  separable little H\"older products. Young local dependence
  and positive endpoint norms make the displayed logarithms continuous;
  the rate is continuous by Proposition~\ref{P:uniform-rate-averaging}.
  Growth bounds make each control-ball supremum finite.
  A countable dense subset of each closed Hilbert ball has the
  same supremum. Supremum over $R$ is an ordinary deterministic
  real supremum, not a random measurability or compact-family claim.
  This proves Corollary~\ref{C:uniform-polynomial-balls}.
\end{proof}

\subsection{Arbitrary small residual paths and the ordered product}

\begin{proposition}[Residual comparison and the integrated rate]
  \label{P:uniform-arbitrary-residual}
  Let $x_m$ be arbitrary zero-start little $C^\theta$ compact-operator
  paths on $[0,B]$, put $v_m=\|x_m\|_\theta$ and
  $Y=(\delta\sum_m v_m^p)^{1/p}$.
  Define
  \[
    P_m^{\rm pert}=S_E(\ell_m)U_{H_m+x_m}(\ell_m,0),\qquad
    z_{\rm rat}=\prod_m\frac{\|P_m^{\rm pert}\|}{\|P_m^{\rm mean}\|}.
  \]
  For fixed integer $B\ge1$, $M\ge0$ and every $K_{\rm res}\ge1$,
  \begin{align}
    \delta\log\max(1,z_{\rm rat})\le{}&
    C_{K_{\rm res}}[(2/B)^{1-1/p}Y+(2/B)^{1-\nu/p}Y^\nu]\notag\\
    &+c_BD_M(K_{\rm res}^{-r}+K_{\rm res}^{\nu-r}),
    \label{E:uniform-residual-ratio}
  \end{align}
  uniformly in $R\ge1$, $\|\Phi\|\le M$.
  The finite constants may be written
  $C_{K_{\rm res}}=C(\theta,B,K_{\rm res})$ and
  $c_B=2^{\nu-1}C_{\rm gr}B$.
  Here $D_M=C_B^r(C_{\rm mean}M)^r$.
  If $b_\delta\ge0$ is any deterministic function of $\delta$ tending to zero
  at each fixed $B$, $M_B$ is as in
  Corollary~\ref{C:uniform-polynomial-balls}, and
  $\mathcal F=2\delta\log\max(1,\|P_N^{\rm pert}\cdots P_0^{\rm pert}\|)$,
  then
  \begin{equation}\label{E:uniform-residual-integrated-rate}
    \limsup_{\substack{B\to\infty\\B\in\mathbb N}}
    \limsup_{\delta\downarrow0}
    \sup_{\substack{R\ge1,\ \|\Phi\|\le M_B\\Y\le b_\delta}}
    \{\mathcal F-2\operatorname{Rate}_R(\Phi)\}\le0 .
  \end{equation}
  No Gaussian, Cameron or Hamiltonian condition on the residuals is imposed.
\end{proposition}

\begin{proof}
  All endpoint norms are positive by the actual Young and free inverses.
  On every cell with $a_m\le K_{\rm res}$ apply
  Proposition~\ref{P:young-upper-multiplicative} with base $H_m$,
  perturbation $x_m$ and left multiplier $S_E(\ell_m)$.
  It gives logarithmic ratio at most
  $C(\theta,B,K_{\rm res})(v_m+v_m^\nu)$.
  On cells with $a_m>K_{\rm res}$ use
  \eqref{E:uniform-product-growth},
  $(a_m+v_m)^\nu\le2^{\nu-1}(a_m^\nu+v_m^\nu)$
  and the actual mean denominator $\|P_m^{\rm mean}\|\ge1$.
  The logarithmic ratio is at most $c_B(1+a_m^\nu+v_m^\nu)$,
  with $c_B$ independent of $K_{\rm res},R,\delta,\Phi,M$.
  Increase $C_{K_{\rm res}}$ to dominate both good and bad
  residual constants. Summation gives a nonnegative bound
  \[
    \log\max(1,z_{\rm rat})
    \le C_{K_{\rm res}}\sum_m(v_m+v_m^\nu)
    +c_B\sum_{a_m>K_{\rm res}}(1+a_m^\nu).
  \]
  The bad sum times $\delta$ is at most
  $D_M(K_{\rm res}^{-r}+K_{\rm res}^{\nu-r})$.
  Finite power mean at exponents $p\ge1$ and $p/\nu\ge1$,
  with nonnegative inputs $v_m,v_m^\nu$, gives
  $\delta\sum v_m\le(2/B)^{1-1/p}Y$ and
  $\delta\sum v_m^\nu\le(2/B)^{1-\nu/p}Y^\nu$.
  This includes $p=\nu$, zero entries and the true empty block,
  with no singular powers. It proves~\eqref{E:uniform-residual-ratio}.

  The product $A_{\rm prod}=\prod_m\|P_m^{\rm mean}\|$ is at least one.
  Submultiplicativity, without commutation, gives
  \[
    \begin{aligned}
      \max(1,\|P_N^{\rm pert}\cdots P_0^{\rm pert}\|)
      &\le A_{\rm prod}\max(1,z_{\rm rat}),\\
      \mathcal F-2\delta\sum_m\log\|P_m^{\rm mean}\|
      &\le2\delta\log\max(1,z_{\rm rat}).
    \end{aligned}
  \]
  At fixed $B,M,K_{\rm res}$ the residual terms vanish on $Y\le b_\delta$.
  Let $\delta\downarrow0$ first and then $K_{\rm res}\to\infty$.
  Since $r>\nu$, the residual excess has limsup at most zero,
  uniformly in $R,\Phi$.
  Now take $M=M_B$ at each fixed $B$ and combine with
  Corollary~\ref{C:uniform-polynomial-balls}; a supremum of sums
  is no larger than the sum of the suprema.
  The outer integer-$B$ limit proves~\eqref{E:uniform-residual-integrated-rate}.
  The endpoint functional is continuous on all tuples by generic
  Young local dependence, including nonphysical residuals.
  Nothing was multiplied as scalar displacement moments.
  This proves Proposition~\ref{P:uniform-arbitrary-residual}.
\end{proof}

These are uniform deterministic $E$-phase product estimates.
Their Gaussian Rayleigh application, original point transfer and
signed-positive score comparison are proved separately in
Sections~\ref{S:actual-gaussian-tubes}--\ref{S:intrinsic-second-moment}.
